% Owner: boundary_d. Independently refereed by root and limits_ab.
\begin{proposition}[B.4.25 (Proved): a one-eigenvalue Hopf boundary has no proper analytic completion]
\label{push:B425}
Let $\alpha(q)=q^m\kappa(q)$, where $m\ge1$ and
$\kappa$ is a holomorphic unit on a disc, and let $\beta$
be holomorphic with $0<|\beta(0)|<1$. Shrink the disc so
$\beta$ is a unit and $0<|\alpha(q)|<|\beta(q)|<1$ for
$q\ne0$. Consider the diagonal primary Hopf surfaces
\[
 H_q=(\C^2\setminus\{0\})/
       \langle(z,w)\mapsto(\alpha(q)z,\beta(q)w)\rangle.
                                                        \tag{B.4.25.1}
\]
Their holomorphic family over the punctured disc has no
proper complex-analytic completion over the full disc,
even allowing singular or nonreduced central fibres and
a singular total space. This remains true after every
finite base change ramified at zero.

In particular one cannot replace the blown-up Hopf
fibres in B.4.21 by the unblown family with one multiplier
$\kappa(q)q\to0$ and the other tending to a nonzero
contraction. This concerns that specified auxiliary
Hopf family. It is not a nonexistence theorem for a
completion of the actual two-torus family in B.4.19.
\end{proposition}
\begin{proof}
\emph{Anticanonical sections of a diagonal Hopf surface.}
A section of $K_{H_q}^{-1}$ lifts to
$f(z,w)\,\partial_z\wedge\partial_w$ on
$\C^2\setminus\{0\}$. Hartogs extension makes $f$
entire on $\C^2$. Deck equivariance is exactly
\[
 f(\alpha z,\beta w)=\alpha\beta f(z,w).
                                                        \tag{B.4.25.2}
\]
For the Taylor coefficient of $z^i w^j$, with $i,j\ge0$,
this says
\[
 \alpha^{i-1}\beta^{j-1}=1.                             \tag{B.4.25.3}
\]
When $i,j\ge1$, the contraction inequalities force
$i=j=1$. When $j=0$, (B.4.25.3) would give
$\beta=\alpha^{i-1}$, which is impossible under
$|\alpha|<|\beta|<1$ (and also for $i\le1$).
When $i=0$, the only possibility is
\[
 \alpha=\beta^N,\qquad j=N+1,\qquad N\ge2.
                                                        \tag{B.4.25.4}
\]
At a fixed parameter there is at most one such $N$,
since $0<|\beta|<1$. Thus
\[
 h^0(H_q,K_{H_q}^{-1})=
 \begin{cases}
 2,&\alpha(q)=\beta(q)^N\text{ for some }N\ge2,\\
 1,&\text{otherwise}.
 \end{cases}                                           \tag{B.4.25.5}
\]
The sections in the first case are represented by
$zw$ and $w^{N+1}$. The computation uses actual
bivector automorphy, with no choice of Picard convention.

\emph{Infinitely many accumulating resonances.}
Put $b=|\beta(0)|\in(0,1)$ and $k=|\kappa(0)|>0$.
For all sufficiently small $q$ we may arrange
$|\kappa(q)|\ge3k/4$ and
\[
 |\beta(q)|\le b\exp(C|q|)
\]
for some fixed $C$, by a holomorphic logarithm of
$\beta(q)/\beta(0)$. Set
\[
 R_N=\left(\frac{2b^N}{k}\right)^{1/m}.
\]
Then $R_N\to0$ and $N R_N\to0$. On $|q|=R_N$, for
all sufficiently large $N$,
\[
 |q^m\kappa(q)|\ge\frac32 b^N,
 \qquad |\beta(q)^N|\le b^N\exp(CN R_N)<\frac54 b^N.
\]
Rouch\'e's theorem shows that
$q^m\kappa(q)-\beta(q)^N$ has $m$ zeros counted with
multiplicity in $|q|<R_N$. None is zero. Choose one,
$q_N$. Different $N$ give different points, since a
nonzero contraction cannot have two equal distinct
positive powers. Hence there are infinitely many
resonances $q_N\to0$ with the first value in
(B.4.25.5).

For each $N$ the resonance function is nonzero at
$q=0$, so is not identically zero. Its zeros in the
punctured disc are discrete. The countable union of
these discrete sets cannot contain an open disc.
There are therefore also nonresonant points arbitrarily
near zero, with the second value in (B.4.25.5).

\emph{Coherent direct images forbid this pattern.}
Suppose first that a proper completion
$\pi:\mathcal X\to\Delta$ has a smooth connected
total threefold and agrees with (B.4.25.1) over
$\Delta^*$. Put $L=K_{\mathcal X}^{-1}$.
On each nonzero fibre, adjunction identifies
$L|_{\mathcal X_q}$ with $K_{H_q}^{-1}$, because
the normal direction is the trivial base direction.
Grauert's proper direct-image theorem gives coherent
sheaves
\[
 M=\pi_*L,\qquad N=R^1\pi_*L
\]
on the disc
\citep[Section 6, Hauptsatz I, p.~59]{Grauert60push}.
Multiplication by a nonzero base function is injective
on $L$, and hence on $M$, since the total space is
smooth and the map is nonconstant. Thus $M$ is
locally free of a fixed rank on the disc. The torsion
subsheaf of $N$ is coherent with locally finite
support. After shrinking, it has no support in
$\Delta^*$.

For any $a\in\Delta^*$, apply the direct-image long
exact sequence to
\[
 0\longrightarrow L\xrightarrow{\ q-a\ }L
 \longrightarrow L|_{\mathcal X_a}\longrightarrow0.
\]
At $a$ its beginning yields
\[
 0\longrightarrow M_a/(q-a)M_a
 \longrightarrow H^0(\mathcal X_a,L|_{\mathcal X_a})
 \longrightarrow \ker(q-a:N_a\to N_a)
 \longrightarrow0.                                   \tag{B.4.25.6}
\]
The last term vanishes. Therefore
$h^0(H_a,K_{H_a}^{-1})=\operatorname{rank}M$ for
\emph{every} sufficiently small $a\ne0$. This
contradicts (B.4.25.5) and the accumulating resonances.
This argument uses coherence and a direct exact
sequence; it does not use the corrected closedness
assertion discussed in the Grauert source audit.

Now suppose an arbitrary proper analytic completion
$\mathcal Z\to\Delta$ exists. Pass to its reduction
and retain the unique irreducible component containing
the connected smooth punctured total space. This is a
closed analytic component, remains proper over the
disc, and dominates the disc. Its restriction over
$\Delta^*$ is exactly the given family. Resolve this
component by a proper analytic resolution chosen to
be an isomorphism on its already nonsingular locus;
see \citet[Theorem 1.6 and the noncompact analytic
extension immediately following it]{BM95AB}.
All singular points lie over zero, so the resolution
is an isomorphism over the entire punctured family.
Properness of the composition is preserved. Near the
compact central fibre, local finiteness of the
analytic resolution permits one shrinking of the
disc; its dominating smooth component is a smooth
proper total threefold as in the preceding paragraph.
Its map to the disc is nonconstant, so the base
function is a non-zero-divisor. No vertical component
was retained which could invalidate that assertion.
The preceding contradiction applies. Reduction also
shows that nilpotent structure supported at zero
cannot evade it.

Finally, a finite base change replaces $\alpha$ by
a holomorphic function of some positive vanishing
order and replaces $\beta$ by a unit whose value
at zero still has modulus strictly between zero and
one. The same proof applies verbatim.
\end{proof}

\begin{remark}[Scope and comparison with the classical toric Hopf degeneration]
There is no conflict with the degeneration in
\citet[Section 13, pp.~116--118]{Oda1978D}.
There the boundary rays are generated by $n$ and
$-n+a n'$ with $a>0$, and the period is translation
by $\lambda^{n'}$. In the universal-cover basis
\[
 v_1=n,\qquad v_2=-n+a n',\qquad
 n'=\frac{v_1+v_2}{a},
\]
the base change $\lambda=t^a$ makes this period
act with eigenvalues $(t,t)$, up to fixed torus
units in the indicated variant. Both eigenvalues
tend to zero. B.4.25 assumes instead a nonzero
limit for the second eigenvalue. The distinction
is forced by the exact period marking of the
candidate Tate interface, and cannot be removed
by citing the existence of some Hopf degeneration.
The blown-up Kato family of B.4.10 is also consistent
with B.4.25: its generic fibres are blow-ups of
$H_q$, so the anticanonical sections must satisfy
the extra vanishing conditions at the marked
blow-up centre and are not the spaces in
(B.4.25.5).
\end{remark}
